\documentclass[11pt]{article}

% BVH Tree for Raytracing

\usepackage[utf8]{inputenc}
\usepackage[T1]{fontenc}
\usepackage{amsmath}
\usepackage{amssymb}
\usepackage{graphicx}
\usepackage{algorithm}
\usepackage{algpseudocode}
\usepackage{booktabs}
\usepackage{hyperref}
\usepackage{geometry}
\geometry{margin=2.5cm}

\newtheorem{invariant}{Invariant}
\newtheorem{lemma}{Lemma}
\newtheorem{theorem}{Theorem}

\title{A Bounding Volume Hierarchy (BVH) for Ray Tracing \\
       \large Structure, construction heuristic, and algorithmic analysis}
\author{Advanced Data Structures Project}
\date{}

\begin{document}
\maketitle

\begin{abstract}
We describe the BVH (Bounding Volume Hierarchy) our ray tracer is built on: what
it partitions, what the decision of \emph{where} to split is based on (the Surface
Area Heuristic, SAH), and why that yields an $\mathcal{O}(\log N)$ query cost. We
derive the SAH from geometric probability, analyze the tree's depth to justify the
cost, and give a correctness argument for the traversal.
\end{abstract}

% ===========================================================================
\section{The problem}

A ray tracer shoots one ray per pixel and looks for the first triangle it hits.
The naive version tests the ray against all $N$ triangles in the scene, so an
image of $P$ pixels costs $\mathcal{O}(P\cdot N)$. With our scene
($N \approx 383{,}000$) that is $\sim 10^{11}$ tests per frame: hopeless. We need a
spatial structure that discards many triangles without testing them.

% ===========================================================================
\section{The structure}

\subsection{Partition objects, not space}

Acceleration structures fall into two families. \textbf{Space subdivision}
(kd-tree, octree, grid) partitions the volume into disjoint regions, and a single
triangle may fall into several. The \textbf{BVH} partitions the \emph{set of
objects}: each triangle goes to \emph{one} leaf, in exchange for the boxes of
sibling nodes being allowed to overlap. Since there is no duplication, memory is
linear and construction is a recursive partition of the triangle array. The
overlap is exactly what the heuristic tries to minimize.

\subsection{The tree and its invariant}

The BVH is a \textbf{binary tree}; each node stores an axis-aligned bounding box
(AABB).

\begin{invariant}\label{inv:bound}
A node's box contains every triangle in its subtree.
\end{invariant}

\noindent
It holds by construction (a node's box is the union of its children's). From it
comes the \textbf{pruning}: if a ray misses a node's box, it misses every triangle
in that subtree and we skip the whole subtree. That is what makes traversal
sublinear. We use AABBs because the ray--box test (\emph{slab} method, Williams et
al.~\cite{williams}) is very cheap and the union of two AABBs is trivial.

\subsection{Memory layout}

A node stores its box and either two child indices (internal) or a range
\texttt{(start, count)} into the triangle array (leaf). Triangles are
\textbf{reordered} during construction so that each leaf's triangles are
contiguous. A tree with $L$ leaves has $2L-1$ nodes and $L \le N$, so there are
$\le 2N-1$ nodes: \textbf{$\mathcal{O}(N)$ memory}.

% ===========================================================================
\section{The heuristic (SAH)}

What decides the quality of the tree is \emph{where} we split at each node. A bad
split leaves big, overlapping boxes that rays enter equally, and little gets
pruned. The \textbf{Surface Area Heuristic} provides the criterion.

\subsection{Cost model and geometric probability}

The expected cost of a node with children $L$ and $R$ is modeled as
\[
C(\text{node}) = C_{\text{trav}} + P(L)\,N_L\,C_{\text{int}} + P(R)\,N_R\,C_{\text{int}},
\]
where $N_L, N_R$ are the triangle counts of each child and $P(\cdot)$ is the
probability that a ray which entered the node also enters that child. By integral
geometry, for a convex $I$ inside a convex $B$ that probability is the ratio of
surface areas,
\[
P(I \mid B) = \frac{S(I)}{S(B)},
\]
which is where the heuristic's name comes from. Normalizing by the parent's area,
the cost of a candidate split is proportional to
\[
\boxed{\;\text{cost} \;\propto\; S(L)\,N_L + S(R)\,N_R\;}
\]
We pick the cheapest split, and compare it against not splitting
($\propto N\,C_{\text{int}}$); if splitting does not help, we make a leaf. That is
what produces leaves holding several triangles.

\subsection{Binning}

Evaluating every possible split would cost $\mathcal{O}(N\log N)$ per node.
Instead (Wald~\cite{wald}) we project the \textbf{centroids} onto the longest axis
of their box --- the axis where they are most spread out, which separates them best
--- divide it into $K{=}12$ \textbf{bins}, and evaluate only the $K-1$ planes
between bins, with area and count accumulated by two linear sweeps. Cost of
choosing the split: $\mathcal{O}(N + K)$ per node.

% ===========================================================================
\section{Algorithms}

\begin{algorithm}[h]
\caption{\textsc{BuildNode}(start, count, depth)}
\begin{algorithmic}[1]
\State \texttt{bounds} $\gets$ union of the boxes of $[start,start+count)$
\If{$count \le 2$ \textbf{or} $depth \ge \text{MAX\_DEPTH}$}
    \State leaf with range $[start, start+count)$; \Return \Comment{base case}
\EndIf
\State $axis \gets$ longest side of the centroid box
\If{centroids nearly coincide on $axis$} leaf; \Return \EndIf
\State bin into $K{=}12$ bins by centroid on $axis$; left/right sweep
\State $i^\star \gets \arg\min_i\; S_L(i)\,n_L(i) + S_R(i)\,n_R(i)$
\If{$1 + \text{cost}(i^\star)/S(\texttt{bounds}) \ge count$} leaf; \Return \Comment{not worth it} \EndIf
\State $mid \gets \textsc{Partition}(start, count, axis, i^\star)$
\If{$mid \in \{start,\, start+count\}$} $mid \gets start + count/2$ \Comment{fallback: median} \EndIf
\State \textsc{BuildNode}$(start, mid-start, depth+1)$
\State \textsc{BuildNode}$(mid, start+count-mid, depth+1)$
\end{algorithmic}
\end{algorithm}

\noindent
\textsc{Partition} reorders the range \emph{in place} with two indices (quicksort
partition scheme), $\mathcal{O}(count)$.

\begin{algorithm}[h]
\caption{\textsc{Intersect}(origin, dir)}
\begin{algorithmic}[1]
\State $best \gets +\infty$;\quad $stack \gets [\,\text{root}\,]$
\While{$stack \neq \emptyset$}
    \State $n \gets$ pop
    \If{the ray misses \texttt{n.bounds} within distance $best$} \textbf{continue} \Comment{prune} \EndIf
    \If{$n$ is a leaf}
        \ForAll{triangle $t \in n$}
            \State $d \gets$ \textsc{MollerTrumbore}(origin, dir, $t$)
            \If{$d$ valid \textbf{and} $d < best$} $best \gets d$; remember $t$ \EndIf
        \EndFor
    \Else\ push both children
    \EndIf
\EndWhile
\State \Return nearest triangle (or ``miss'')
\end{algorithmic}
\end{algorithm}

\noindent
Shadows use \textsc{Occluded}: identical, but it returns \emph{true} at the first
triangle between the point and the light (no need for the nearest one).

% ===========================================================================
\section{Depth and split analysis}

\subsection{Construction}

The work at a node (boxes, binning, partition) is $\Theta(N)$, so
\[
B(N) = \Theta(N) + B(N_L) + B(N_R), \qquad N_L + N_R = N.
\]
With balanced splits ($N_L, N_R \approx N/2$), by the Master Theorem
$B(N) = \Theta(N \log N)$: each level distributes the $N$ triangles without
repetition ($\Theta(N)$) and there are $\Theta(\log N)$ levels.

\subsection{Why are the splits balanced?}

The SAH does not optimize balance directly, but it favors it: a very unbalanced
split ($N_L \approx N$) leaves a box almost as large as the parent's holding almost
all the triangles, so its term $S(L)\,N_L$ barely drops versus not splitting and
the heuristic rejects it. With reasonably uniform geometry this gives, level by
level, a geometric reduction $N \to \alpha N$ ($\alpha < 1$ bounded away from $1$),
i.e.\ expected depth $\mathcal{O}(\log N)$; in practice SAH BVHs measure heights
close to $\log_2 N$. Moreover, when binning degenerates the algorithm falls back to
a \textbf{median} split, which halves the set and \emph{guarantees}
$\le \lceil \log_2 N\rceil$ at that point.

\subsection{Worst case}

If many triangles share a centroid, no plane separates them and the partition may
peel one off per level: a chain of depth $\Theta(N)$. We bound it by making a leaf
when the centroid box is degenerate, and with a cap \texttt{MAX\_DEPTH}$=60$ that
forces a leaf upon reaching it. This bounds the height and, with it, the traversal
stack (fixed, 64), at the cost of a slightly fatter leaf here and there.

\subsection{Query}

Under the SAH model (small, weakly overlapping child boxes) a ray visits
$\mathcal{O}(\log N)$ nodes plus a constant number of triangles per leaf; the
$best$ clamp trims further. The worst case (a ray grazing many overlapping boxes,
or a degenerate tree) is $\mathcal{O}(N)$, rare with a SAH BVH.

\begin{table}[h]
\centering
\begin{tabular}{lll}
\toprule
Operation & Average & Worst case \\
\midrule
Construction & $\mathcal{O}(N \log N)$ & $\mathcal{O}(N^2)$ (bounded by MAX\_DEPTH) \\
Memory & $\mathcal{O}(N)$ & $\mathcal{O}(N)$ \\
Nearest hit / shadow & $\mathcal{O}(\log N)$ & $\mathcal{O}(N)$ \\
Frame ($P$ pixels) & $\mathcal{O}(P \log N)$ & $\mathcal{O}(P\,N)$ \\
\bottomrule
\end{tabular}
\caption{Cost of each operation.}
\end{table}

% ===========================================================================
\section{Correctness}

\begin{lemma}[Partition]
The leaves form a partition of the triangles: each one is in exactly one leaf.
\end{lemma}
\noindent
\textsc{Partition} returns $mid$ with $[start,mid)$ and $[mid,start+count)$
disjoint and covering the range; the recursion descends into those subranges and
the base case assigns a full range to a leaf. By induction the leaves' ranges are
disjoint and cover $[0,N)$. Since $count$ strictly decreases (or is cut by
\texttt{MAX\_DEPTH}), the recursion \textbf{terminates}.

\begin{lemma}[Sound pruning]\label{lem:prune}
If the ray misses \texttt{n.bounds}, it misses every triangle in $n$'s subtree.
\end{lemma}
\noindent
Immediate from Invariant~\ref{inv:bound}.

\begin{theorem}[Correctness of \textsc{Intersect}]
It returns the triangle of smallest $t>0$ along the ray, the same as testing all
$N$.
\end{theorem}
\noindent
A node is discarded only if (i) the ray misses its box (Lemma~\ref{lem:prune}), or
(ii) it is hit at distance $\ge best$, in which case by Invariant~\ref{inv:bound}
every triangle in the subtree has $t \ge best$ and cannot improve the optimum.
Every node not discarded is visited, and at each leaf all its triangles are tested,
so the final $best$ is the global minimum. Each node is pushed once over a finite
tree, hence it \textbf{terminates}. \textsc{Occluded} uses the same pruning with
the weaker condition ``some $t < t_{\text{light}}$''.

% ===========================================================================
\section{Integration in the ray tracer}

On top of the BVH: a \textbf{static} hierarchy for the room (built once) and a
\textbf{dynamic} one rebuilt per frame for the animated shapes; each ray queries
both. Shadows via \textsc{Occluded}, recursive Whitted-style reflections
($\text{color} = \text{local}(1-k) + \text{reflected}\cdot k$), ambient lighting
from a \emph{cubemap} (skybox) sampled along the normal, and parallelism by
splitting pixels across threads (BVH queries are read-only).

% ===========================================================================
\section{Results}

\begin{figure}[h]
  \centering
  \fbox{\parbox[c][4cm][c]{0.46\textwidth}{\centering FIGURE: rasterized
  (docs/img/raster.png)}}
  \hfill
  \fbox{\parbox[c][4cm][c]{0.46\textwidth}{\centering FIGURE: raytrace + BVH
  (docs/img/raytrace.png)}}
  \caption{Same scene rasterized and ray traced over the BVH.}
\end{figure}

\noindent
\emph{(Fill in your numbers: $N$, FPS with and without the BVH, thread count, build
time, measured tree height versus $\log_2 N$, nodes visited per ray.)}

% ===========================================================================
\begin{thebibliography}{9}

\bibitem{goldsmith}
J. Goldsmith, J. Salmon.
\emph{Automatic Creation of Object Hierarchies for Ray Tracing}.
IEEE Computer Graphics and Applications, 7(5), 1987.

\bibitem{macdonald}
D. J. MacDonald, K. S. Booth.
\emph{Heuristics for Ray Tracing Using Space Subdivision}.
The Visual Computer, 6(3), 1990.

\bibitem{wald}
I. Wald.
\emph{On Fast Construction of SAH-based Bounding Volume Hierarchies}.
IEEE Symposium on Interactive Ray Tracing, 2007.

\bibitem{moller}
T. M\"oller, B. Trumbore.
\emph{Fast, Minimum Storage Ray/Triangle Intersection}.
Journal of Graphics Tools, 2(1), 1997.

\bibitem{williams}
A. Williams, S. Barrs, R. K. Morley, P. Shirley.
\emph{An Efficient and Robust Ray-Box Intersection Algorithm}.
Journal of Graphics Tools, 2005.

\bibitem{pbrt}
M. Pharr, W. Jakob, G. Humphreys.
\emph{Physically Based Rendering: From Theory to Implementation} (4th ed.).
MIT Press, 2023. Chapter 4.

\end{thebibliography}

\end{document}
